% ════════════════════════════════════════════════════════════
\section{CAPÍTULO II: MARCO TEÓRICO Y CONCEPTUAL}
% ════════════════════════════════════════════════════════════

\subsection{Antecedentes}

El estudio de los ecosistemas de paquetes de software como redes complejas 
ha ganado atención creciente en la comunidad de ingeniería de software empírica.
Los trabajos existentes pueden agruparse en tres líneas principales.

\textbf{Estructura topológica de grafos de dependencias.}
\textcite{kikas2017structure} analizaron los grafos de dependencias de npm, 
PyPI y RubyGems, encontrando que los tres presentan distribución de grado 
con cola pesada, característica de redes libres de escala, con un número 
reducido de paquetes concentrando la mayoría de las dependencias entrantes. 
Sin embargo, su trabajo se limitó a métricas topológicas descriptivas 
(densidad, diámetro, coeficiente de agrupamiento) sin proponer un índice 
de criticidad sistémica para los nodos individuales.

\textbf{Propagación de vulnerabilidades en ecosistemas de software.}
\textcite{decan2019impact} estudiaron el impacto transitivo de vulnerabilidades 
de seguridad registradas en el ecosistema npm, demostrando que un subconjunto 
pequeño de paquetes concentra una proporción desproporcionada del riesgo 
propagado. No obstante, su enfoque es retrospectivo: parte de CVEs ya 
conocidos para medir el daño posterior, sin proponer una métrica proactiva 
capaz de anticipar dicho riesgo antes del incidente. De forma complementaria, 
\textcite{zimmermann2019small} identificaron que más del 40\% de los paquetes 
npm dependen transitivamente de al menos un paquete con un único mantenedor, 
señalando la fragilidad humana del ecosistema, pero sin formalizar la 
propagación mediante un modelo estocástico.

\textbf{Aplicaciones de PageRank y centralidad en redes de software.}
El algoritmo PageRank, introducido por \textcite{brin1998anatomy} como 
distribución estacionaria de una Cadena de Markov sobre el grafo web, ha 
sido aplicado posteriormente a redes de citas académicas, redes de 
colaboración en proyectos open source y grafos de llamadas entre funciones. 
Sin embargo, su aplicación directa al grafo de dependencias de PyPI como 
índice de riesgo sistémico no ha sido explorada en la literatura revisada, 
constituyendo el aporte central del presente trabajo. El algoritmo de 
detección de comunidades de \textcite{blondel2008louvain} ha sido empleado 
en redes de co-autoría y co-contribución en GitHub, pero tampoco se ha 
aplicado al análisis espacial del riesgo en ecosistemas de dependencias.

En síntesis, los antecedentes identifican la existencia del problema 
(concentración de riesgo en ecosistemas de paquetes) y caracterizan 
su estructura topológica, pero ninguno propone un modelo estocástico 
formal que combine Cadenas de Markov, PageRank y análisis de comunidades 
para construir un índice de riesgo sistémico proactivo sobre PyPI.


\subsection{Marco Teórico}

\subsubsection{El Grafo de Dependencias como Espacio de Probabilidad}
Sea $\Omega = V = \{p_1, p_2, \ldots, p_n\}$ el conjunto de paquetes de PyPI
(nodos del grafo), $\mathcal{F} = 2^\Omega$ la $\sigma$-álgebra potencia, y
$\mathbb{P}$ una medida de probabilidad sobre $(\Omega, \mathcal{F})$. La
distribución estacionaria PageRank $\mathbf{r}^*$ es precisamente la medida
$\mathbb{P}$ invariante de la Cadena de Markov definida sobre este espacio.

La relación de dependencia define un grafo dirigido $G = (V, E)$ donde la arista
$(p_i, p_j) \in E$ indica que el paquete $p_i$ \textit{depende de} $p_j$.
El \textbf{random walk sobre dependencias} modela a un agente que, dado que se
encuentra en el paquete $p_i$, salta uniformemente a cualquiera de las
dependencias de $p_i$.

\subsubsection{Cadena de Markov sobre el Grafo de Dependencias}
\begin{definicion}[Matriz de hipervínculos de dependencias $\mathbf{H}$]
Sea $\mathbf{L} \in \{0,1\}^{n \times n}$ la matriz de adyacencia del grafo
dirigido, con $L_{ij}=1$ si $p_j$ depende de $p_i$. La \textbf{matriz de
transición de dependencias} es:
\[
    H_{ij} =
    \begin{cases}
        \dfrac{L_{ij}}{\text{outdeg}(j)} & \text{si } \text{outdeg}(j) > 0 \\[6pt]
        0 & \text{si } \text{outdeg}(j) = 0
    \end{cases}
\]
$\mathbf{H}$ es sub-estocástica por columnas: los paquetes sin dependencias
salientes (librerías base como \texttt{numpy}) generan columnas nulas.
\end{definicion}

\begin{definicion}[Corrección de nodos sumidero]
Sea $\mathbf{a} \in \{0,1\}^n$ el vector indicador de paquetes sin dependencias
salientes ($a_j = 1 \Leftrightarrow \text{outdeg}(j)=0$). La \textbf{matriz
estocástica corregida} es:
\[
    \mathbf{A} = \mathbf{H} + \frac{1}{n}\mathbf{1}\mathbf{a}^\top
\]
Intuitivamente: cuando el agente llega a un paquete sin dependencias, se
\textit{teleporta} uniformemente a cualquier paquete del ecosistema.
$\mathbf{A}$ es estocástica por columnas: $\mathbf{1}^\top \mathbf{A} =
\mathbf{1}^\top$.
\end{definicion}

\begin{definicion}[Matriz de Google aplicada a dependencias $\mathbf{G}(d)$]
Para $d \in (0,1)$ (factor de amortiguamiento), la \textbf{matriz de Google
del ecosistema} es:
\[
    \boxed{\mathbf{G}(d) = d\,\mathbf{A} + \frac{1-d}{n}\mathbf{E}}
\]
donde $\mathbf{E} = \mathbf{1}\mathbf{1}^\top$. Con probabilidad $d$ el agente
sigue una dependencia real; con probabilidad $1-d$ salta a un paquete aleatorio
del ecosistema. $\mathbf{G}(d)$ es positiva, irreducible y aperiódica, lo que
garantiza, por el Teorema de Perron–Frobenius, la existencia y unicidad de
$\mathbf{r}^*$.
\end{definicion}

\begin{teorema}[Perron–Frobenius aplicado al ecosistema PyPI]
Sea $\mathbf{G}(d)$ la matriz de Google del ecosistema de dependencias con
$d \in (0,1)$. Entonces:
\begin{enumerate}
    \item $\lambda_1 = 1$ es el autovalor dominante simple de $\mathbf{G}(d)$.
    \item El autovector correspondiente $\mathbf{r}^*$ tiene componentes
          estrictamente positivas y $\|\mathbf{r}^*\|_1 = 1$.
    \item Para cualquier distribución inicial $\mathbf{r}^{(0)}$:
          \[
              \lim_{k\to\infty} \mathbf{G}(d)^k \mathbf{r}^{(0)} = \mathbf{r}^*,
              \quad \text{con tasa } \|\mathbf{r}^{(k)} - \mathbf{r}^*\|_1 = O(d^k).
          \]
\end{enumerate}
\end{teorema}

\subsubsection{Interpretación del PageRank como Riesgo Sistémico}
El PageRank $r_i^*$ de un paquete $p_i$ puede interpretarse de dos formas
equivalentes y complementarias:
\begin{enumerate}
    \item \textbf{Probabilística}: es la fracción de tiempo que un agente que
          navega aleatoriamente el grafo de dependencias pasa en el paquete $p_i$.
          Un valor alto indica que $p_i$ es ``ineludible'' en el ecosistema.
    \item \textbf{Estructural}: es proporcional a la suma ponderada de los
          PageRanks de todos los paquetes que dependen de $p_i$, de forma
          transitiva. Captura tanto la centralidad directa (muchos usuarios)
          como la transitiva (usuarios de usuarios).
\end{enumerate}
Esta segunda interpretación es la que distingue al PageRank del simple in-degree:
un paquete usado por pocos paquetes de altísima centralidad puede tener más
PageRank que uno usado directamente por muchos paquetes de baja centralidad.

\subsubsection{Método de la Potencia (Power Iteration)}
El algoritmo iterativo aplica repetidamente:
\[
    \mathbf{r}^{(k+1)} = \mathbf{G}(d)\,\mathbf{r}^{(k)}
    = d\!\left(\mathbf{H}\mathbf{r}^{(k)} +
      \frac{\mathbf{a}^\top\mathbf{r}^{(k)}}{n}\mathbf{1}\right)
      + \frac{1-d}{n}\mathbf{1}
\]
La formulación explota la estructura dispersa de $\mathbf{H}$ (el grafo de
PyPI tiene densidad $\delta \ll 1$) y evita materializar $\mathbf{G}(d)$
en memoria, lo que es crítico para redes de cientos de miles de nodos.

\subsubsection{Detección de Comunidades y Análisis Espacial}
El análisis espacial del ecosistema se realiza mediante la detección de
comunidades con el \textbf{algoritmo de Louvain}, que maximiza la modularidad:
\[
    Q = \frac{1}{2m} \sum_{ij} \left[ A_{ij} - \frac{k_i k_j}{2m} \right]
    \delta(c_i, c_j)
\]
donde $A_{ij}$ es la matriz de adyacencia, $k_i$ el grado del nodo $i$,
$m$ el número total de aristas y $\delta(c_i, c_j) = 1$ si los nodos $i$
y $j$ pertenecen a la misma comunidad. Los clusters resultantes representan
subcomunidades temáticas del ecosistema (data science, web, criptografía, etc.)
que pueden visualizarse como un \textbf{mapa espacial del riesgo}.

\subsection{Marco Conceptual}

Los siguientes conceptos indican el vocabulario técnico del presente trabajo.

\textbf{Ecosistema de software.}
Conjunto de paquetes de software publicados en un gestor de paquetes (en este 
caso PyPI), junto con las relaciones de dependencia entre ellos, sus 
mantenedores y su comunidad de usuarios. El término ``ecosistema'' enfatiza 
que los paquetes no existen de forma aislada sino en una red de 
interdependencias funcionales \autocite{kikas2017structure}.

\textbf{Grafo dirigido de dependencias.}
Representación formal $G = (V, E)$ donde cada nodo $v \in V$ es un paquete 
de PyPI y cada arista dirigida $(p_i, p_j) \in E$ indica que el paquete 
$p_i$ declara al paquete $p_j$ como dependencia \texttt{runtime}. La 
dirección de la arista modela la relación de dependencia: $p_i$ requiere 
a $p_j$ para funcionar.

\textbf{Dependencia transitiva.}
Si el paquete $A$ depende de $B$, y $B$ depende de $C$, entonces $A$ 
depende transitivamente de $C$, aunque no lo declare explícitamente. La 
vulnerabilidad o fallo de $C$ afecta a $A$ a través de la cadena 
$A \to B \to C$. Las dependencias transitivas son invisibles para análisis 
locales pero capturadas por el random walk de la Cadena de Markov.

\textbf{Riesgo sistémico.}
En el contexto de ecosistemas de software, el riesgo sistémico de un paquete 
es la magnitud del daño en cascada que su fallo o compromiso generaría sobre 
el resto del ecosistema. Se distingue del riesgo local (si el paquete tiene 
una vulnerabilidad activa) por ser una propiedad estructural de la red, 
independiente de la existencia previa de CVEs registrados.

\textbf{Nodo sumidero (\textit{dangling node}).}
Paquete sin dependencias salientes declaradas (\texttt{out-degree} $= 0$). 
Típicamente corresponde a librerías base altamente estables como 
\texttt{numpy} o \texttt{cryptography}. Su presencia genera columnas nulas 
en la matriz de transición $\mathbf{H}$, rompiendo la condición de 
estocasticidad y requiriendo corrección explícita mediante la matriz 
$\mathbf{A}$.

\textbf{Cadena de Markov ergódica.}
Proceso estocástico en tiempo discreto $\{X_t\}_{t \geq 0}$ sobre un espacio 
de estados finito $V$, que satisface la propiedad de Markov 
($P(X_{t+1} | X_t, \ldots, X_0) = P(X_{t+1} | X_t)$) y es ergódico 
(irreducible y aperiódico). La ergodicidad garantiza la existencia y unicidad 
de una distribución estacionaria $\boldsymbol{\pi}$ tal que 
$\boldsymbol{\pi} = \mathbf{P}\boldsymbol{\pi}$.

\textbf{Distribución estacionaria (PageRank).}
Vector de probabilidad $\mathbf{r}^* \in \mathbb{R}^n$ invariante bajo la 
acción de la matriz de transición: $\mathbf{r}^* = \mathbf{G}(d)\,\mathbf{r}^*$. 
Representa la fracción de tiempo que un agente que navega el grafo de 
dependencias indefinidamente pasa en cada paquete. En este trabajo, 
$\mathbf{r}^*$ constituye el índice de riesgo sistémico.

\textbf{Factor de amortiguamiento ($d$).}
Parámetro $d \in (0,1)$ de la matriz de Google $\mathbf{G}(d)$ que modela 
la probabilidad de que el agente siga una dependencia real en lugar de 
teleportarse a un paquete aleatorio. Controla el balance entre centralidad 
local (bajo $d$) y centralidad global transitiva (alto $d$), y determina 
la tasa de convergencia geométrica $O(d^k)$ del método de la potencia.

\textbf{Centralidad local vs.\ centralidad transitiva.}
La centralidad local de un paquete es su in-degree: cuántos otros paquetes 
lo usan directamente. La centralidad transitiva (PageRank) pondera además 
la importancia de quienes lo usan: un paquete usado por pocas librerías 
muy importantes puede tener mayor PageRank que uno usado por muchos paquetes 
de baja relevancia. Esta distinción es el fundamento de la hipótesis 
principal del trabajo.

\textbf{Comunidad o \textit{cluster} del ecosistema.}
Subconjunto de nodos del grafo con mayor densidad de aristas entre sí que 
hacia el resto de la red. En el ecosistema PyPI, las comunidades corresponden 
a subcomunidades temáticas (data science, desarrollo web, criptografía, etc.). 
Su detección mediante el algoritmo de Louvain permite analizar la distribución 
\textit{espacial} del riesgo sistémico entre distintas áreas funcionales del ecosistema.
